% Generated draft: requires mathematical and source review.
\documentclass[11pt]{article}
\usepackage[margin=25mm]{geometry}
\usepackage{fontspec}
\setmainfont{Latin Modern Roman}
\usepackage{amsmath,amssymb,mathtools,unicode-math}
\setmathfont{Latin Modern Math}
\usepackage{enumitem,needspace,xurl,hyperref}
\setlength{\parindent}{0pt}
\setlength{\parskip}{5pt}
\setlength{\emergencystretch}{3em}
\newcommand{\N}{\mathbb{N}}
\newcommand{\Z}{\mathbb{Z}}
\newcommand{\Q}{\mathbb{Q}}
\newcommand{\R}{\mathbb{R}}
\newcommand{\C}{\mathbb{C}}
\newcommand{\F}{\mathbb{F}}
\newcommand{\A}{\mathbb{A}}
\newcommand{\PP}{\mathbb{P}}
\newcommand{\E}{\mathbb{E}}
\newcommand{\eps}{\varepsilon}
\newcommand{\dd}{\,\mathrm{d}}
\begin{document}
\section*{UnsolvedMath ID: 30004063}
\textbf{Source problem code:} catalogue entry
\subsection{Definitions and mathematical statement}
Let \(\mathcal X\) denote the admissible class of binary images and let \(\mathcal T\) denote the tomographic forward map for a given instance. For noiseless data \(p\), define the feasible set by
\[
\mathcal S(\mathcal T,p)=\{x\in\mathcal X:\mathcal T(x)=p\}.
\]
The instance has a unique binary solution if \(\mathcal S(\mathcal T,p)=\{x_\ast\}\) for some \(x_\ast\in\mathcal X\).

The supplied question asks whether, for every instance in the intended class, uniqueness of the binary solution guarantees that the proposed Lagrangian-dual approach retrieves \(x_\ast\). If that approach defines a reconstruction map \(\mathcal D\), the requested implication is schematically
\[
\mathcal S(\mathcal T,p)=\{x_\ast\}
\quad\Longrightarrow\quad
\mathcal D(\mathcal T,p)=x_\ast.
\]
Here \(\mathcal X\), \(\mathcal T\), and \(\mathcal D\) are placeholders for the intended image class, forward model, and reconstruction procedure; their definitions are not supplied. In particular, the record does not specify the binary values, image domain, acquisition geometry, boundary conditions, admissible operators, dual optimization problem, or reconstruction rule. Thus the displayed implication records the intended recovery question but is not a fully specified conjecture.

The supplied statement-verification note classifies this record as \texttt{unrecoverable} and describes it as a partial duplicate of the binary-tomography recovery conjecture retained under source code \texttt{OWR-16766-003}. It explicitly declines to promote a separate clean statement for this record. That related conjecture is not supplied here and cannot be substituted for this question.
\subsection{Short English statement}
For noiseless binary tomography, does a unique feasible binary image necessarily coincide with the reconstruction produced by the proposed Lagrangian-dual method? The method and its admissible instances are unspecified. The supplied registry retains this record for provenance as a partial duplicate, so this is a schematic statement requiring human mathematical review.
\subsection{Sources}
\begin{itemize}
\item Local UnsolvedMath registry, supplied complete record: permanent ID \(30004063\), title \emph{Exact Recovery by Lagrangian-Dual Binary Tomography}, exact source code \texttt{catalogue entry}, collection \emph{Oberwolfach Reports Open Problems}. The recorded proposed year is 2019; no proposer is identified.
\item Supplied report citation: \emph{Tomographic Inverse Problems: Theory and Applications} (2020), Oberwolfach Reports, pp.~209--303, DOI \texttt{10.4171/owr/2019/4}. \url{https://doi.org/10.4171/owr/2019/4}. The background describes the reports as documenting mathematical workshops and their open-problem discussions.
\item Supplied literature source: Kadu--van Leeuwen, convex binary tomography. \url{https://doi.org/10.1109/TCI.2019.2898333}.
\item Literature assessment dated exactly \texttt{2026-08-22}: the supplied curation labels the problem open, describes support for the dual method as limited to exhaustive small-image experiments, and reports that the primary paper gives no general guarantee of recovery of an optimal binary solution and does not prove the conjectured uniqueness or intersection property. This is supplied dated curation derived from the OWR record, not an independently verified literature assessment; its status and cited claims require independent verification.
\item Supplied statement-verification note: the official report is described as showing a partial duplicate of the conjecture retained in \texttt{OWR-16766-003}; this record remains visible for provenance, with no separate clean statement promoted. This source-verification claim is reproduced from the record and has not been independently checked.
\end{itemize}
\end{document}
